\subsection{Instruction Payload Ordering}
The instruction header is byte 0 for extrashort and short instructions and the first two bytes for medium and
longer instructions. The opcode space may continue after the header. If a descriptor, immediate, displacement, or
bitmap follows the opcode space, each independently determined unit is a separate payload, and the payloads occupy
increasing addresses in decode order. An extended EA descriptor and a following displacement are therefore separate
payloads.

Signed displacement and immediate payloads use two's-complement representation at their declared width before any
sign extension performed by the consuming instruction or effective-address encoding.

\Needspace{1.25in}
\manualtablecaption{Immediate Operand Interpretation}
\begin{manuallongtable}{@{}>{\raggedright\arraybackslash}p{0.95in}>{\raggedright\arraybackslash}p{0.38in}>{\raggedright\arraybackslash}p{0.62in}>{\raggedright\arraybackslash}p{0.62in}>{\raggedright\arraybackslash}p{0.90in}>{\raggedright\arraybackslash}p{1.70in}@{}}
\toprule
\rowcolor{ManualHeaderFill}
\textbf{Operand} & \textbf{Bits} & \textbf{Value} & \textbf{Range} & \textbf{Extension} & \textbf{Applies When}\\
\midrule
\endfirsthead
\multicolumn{6}{@{}l}{\scriptsize\itshape Table \themanualtable\ (continued)}\\
\toprule
\rowcolor{ManualHeaderFill}
\textbf{Operand} & \textbf{Bits} & \textbf{Value} & \textbf{Range} & \textbf{Extension} & \textbf{Applies When}\\
\midrule
\endhead
\texttt{PAIRn} & 3 & identifier & 0..7 & none & PUSHP/POPP canonical register-pair selector\\
\texttt{pt\_level} & 3 & enumeration & 1..5 & none & PTQUERY page-table level selector\\
\texttt{flags\_bitmap} & 4 & bitmap & 0..15 & none & SETF integer FLAGS image\\
\texttt{imm6} & 6 & unsigned & 0..63 & zero-extend & integer encodings with an explicit B/W/L/Q operation size\\
\texttt{imm7} & 7 & unsigned & 0..127 & zero-extend & EXTRACT register-pair encodings\\
\bottomrule
\end{manuallongtable}


PUSHP and POPP interpret the @PAIR_ID_WIDTH@-bit \texttt{pair\_id} as a canonical register-pair index, so every encoding
from @PAIR_ID_RANGE@ is valid.
PTQUERY uses a @PT_LEVEL_WIDTH@-bit \texttt{pt\_level}; assigned levels are @PT_LEVEL_RANGE@, while encodings 0, 6, and 7
are reserved.
SETF interprets the @FLAGS_BITMAP_WIDTH@-bit \texttt{flags\_bitmap} range @FLAGS_BITMAP_RANGE@ as a complete FLAGS
image: bit 3 writes Z, bit 2 writes N, bit 1 writes C, and bit 0 writes V. Each one writes the corresponding flag to
one, and each zero writes it to zero.

\begin{manuallistedformatdiagram}{SETF Flag Bitmap}
\manualformatrowrange{flags\_bitmap[3:0]}{3}{0}{%
\manualformatfield{Z}{1}
\manualformatfield{N}{1}
\manualformatfield{C}{1}
\manualformatfield{V}{1}
}
\end{manuallistedformatdiagram}

A @IMM6_WIDTH@-bit \texttt{imm6} value in @IMM6_RANGE@ is zero-extended to the selected operation size before use when
that size is wider than @IMM6_EXTENSION_THRESHOLD@ bits. EXTRACT uses the @IMM7_WIDTH@-bit \texttt{imm7} range @IMM7_RANGE@ as the
offset into its concatenated register pair.
